\section{Support-stable local-linearization diagnostic}
\label{sec:cstab_local_linearizations}

The main text reports the simpler \(F_{\rm abs}\)-routed local affine predictor, because matched route controls show that instantaneous support families are already sufficient to improve over the global-\(K\) LISTA baseline and are comparable to, or slightly better than, the more elaborate stable support component route. We keep the \(C_{\rm stab}\) construction as an appendix diagnostic because it asks a different mechanistic question: whether support families induced by the sparse encoder have a coherent support-flow fate.

\paragraph{Stable support components.}
The object \(C_{\rm stab}\) starts from the same absolute mask \(m_{\rm abs}(x)=\mathbf{1}\{|z_i|>10^{-3}\}\), but first forms a finer base support-family label \(u_t\) using Jaccard threshold \(0.8\). On route-fitting trajectories, consecutive labels define an empirical support-transition graph with edges \(u\to v\). We filter this graph by transition frequency, compute strongly connected components, retain recurrent components that appear in trajectory tails and have small outbound probability, and assign each observed base label \(u\) to the recurrent component into which its future support trajectory is absorbed with high empirical confidence. Thus \(C_{\rm stab}(u)\) is a support-flow fate label, not a basin label, fixed-point estimate, or prescribed switching mode. Basin labels and basin counts are not used to construct \(C_{\rm stab}\); they are used only afterward for benchmark audits.

\begin{figure}[tbp]
\centering
\resizebox{0.98\linewidth}{!}{%
\begin{tikzpicture}[
  font=\scriptsize,
  box/.style={draw,rounded corners=2pt,align=center,inner xsep=5pt,inner ysep=4pt,text width=2.45cm,minimum height=0.92cm},
  route/.style={-{Latex[length=2mm]},thick}
]
  \node[box] (x) {state\\ \(x_t\)};
  \node[box,right=0.45cm of x] (z) {encoder\\ \(z_t=\Enc(x_t)\)};
  \node[box,right=0.45cm of z,text width=3.0cm] (mask) {absolute mask\\ \(m_{\rm abs}(x_t)\)};
  \node[box,right=0.45cm of mask,text width=2.75cm] (base) {base support-family\\ label \(u_t\)\\ Jaccard \(0.8\)};
  \node[box,below=0.85cm of mask,text width=3.2cm] (graph) {support-transition graph\\ edge \(u\to v\)};
  \node[box,right=0.55cm of graph,text width=3.3cm] (stab) {recurrent components\\ and absorption fates\\ \(c=C_{\rm stab}(u_t)\)};
  \node[box,right=0.55cm of stab,text width=3.3cm] (local) {route-local affine map\\ \(T_c(z)=d_c+K_c(z-\bar z_c)\)};
  \draw[route] (x) -- (z);
  \draw[route] (z) -- (mask);
  \draw[route] (mask) -- (base);
  \draw[route] (base.south) |- (graph.north);
  \draw[route] (graph) -- (stab);
  \draw[route] (stab) -- (local);
  \node[draw,dashed,rounded corners=2pt,fit=(graph)(stab),inner sep=6pt,label={[font=\scriptsize]below:built once after stage 1 from training trajectories}] {};
\end{tikzpicture}}
\caption{\(C_{\rm stab}\) routing used by the staged local-linearization diagnostic. Exact support masks define base labels \(u_t\), transitions define a support-flow graph, recurrent support-flow components define \(C_{\rm stab}\), and \(c=C_{\rm stab}(u_t)\) selects the local affine latent map.}
\label{fig:cstab_routing}
\end{figure}

\paragraph{Staged training.}
The staged training recipe is the same as the main \(F_{\rm abs}\)-routed local-map experiment. Given total budget \(T=200{,}000\), the first \(100{,}000\) steps train the LISTA Koopman autoencoder with a shared encoder \(\Enc\), decoder \(\Dec\), and global latent matrix \(K\). We then freeze \(\Enc\), \(\Dec\), and \(K\), build \(C_{\rm stab}\) from \(512\) training-split trajectories of length \(192\), and train one affine latent map for each retained component for the final \(100{,}000\) steps. For component \(c\), the predictor is
\begin{equation}
  T_c(z)=d_c+K_c(z-\bar z_c),
  \label{eq:cstab_local_affine}
\end{equation}
where \(\bar z_c\) is the mean source latent assigned to the component. Each \(K_c\) is initialized from the frozen global \(K\), and \(d_c\) is initialized to \(K\bar z_c\) and then learned. If the current support mask cannot be assigned to a trained stable component, the predictor falls back to the frozen global map for that step.

\paragraph{Forecasting result and matched-route interpretation.}
Against the matched same-budget global-\(K\) LISTA baseline, the \(C_{\rm stab}\)-routed local affine maps win \(188/225\) paired rows at \(H1000\), with median staged/global ratio \(0.399\). This confirms that support-derived routes can select useful local affine latent dynamics. However, the matched route-baseline packet shows that this is not specific to \(C_{\rm stab}\): instantaneous \(F_{\rm abs}\) support-family routing wins \(189/225\) rows at \(H1000\), with median ratio \(0.328\), and latent \(k\)-means routing is also comparable. We therefore treat \(C_{\rm stab}\) as an interpretable support-flow diagnostic rather than as the primary forecasting route.

\begin{table}[tbp]
\centering
\caption{\(C_{\rm stab}\)-routed staged local affine forecasting on the controlled multibasin benchmark. The baseline is the matched same-budget global-\(K\) LISTA model, with both models evaluated using the same periodic re-encoding grid. Lower staged/global ratios are better.}
\label{tab:cstab_local_k_forecasting}
{\setlength{\tabcolsep}{4pt}
\resizebox{0.86\textwidth}{!}{\begin{tabular}{lcccc}
\toprule
Metric & H100 & H500 & H1000 & All horizons \\
\midrule
Paired staged wins & $189/225$ & $188/225$ & $188/225$ & $176/225$ \\
Geometric staged/global MSE ratio & $0.269$ & $0.202$ & $0.182$ & -- \\
Systems won by seed-geometric ratio & $14/15$ & $14/15$ & $14/15$ & $14/15$ \\
Two-sided paired sign-test \(p\) & $3.0{\times}10^{-26}$ & $1.6{\times}10^{-25}$ & $1.6{\times}10^{-25}$ & $5.3{\times}10^{-18}$ \\
\bottomrule
\end{tabular}
}}
\end{table}
